\documentclass{article}
	\usepackage[UTF8]{ctex}
	\usepackage{amsmath}
	\usepackage{graphicx}
	\usepackage{underscore}
	\usepackage{geometry}
	\usepackage{anysize}
	\date{}
	\title{奋斗折线 \\ [2ex] \begin{large} 
	时间限制：C/C++ 1秒，其他语言2秒 \\
	空间限制：C/C++ 256.000MB，其他语言512.000MB 	\\
	$64$bit IO Format: \%lld
 			\end{large} }
	\geometry{top=10mm,bottom=15mm,left=25mm,right=25mm}
	
	\begin{document}
		\maketitle \vspace{-50pt}
		
		\section*{题目描述}
		\indent 编程比赛的征程，道路是曲折的。当下素材已经呈上，请为自己送上一条奋斗折线，激励自己勇往前行！\\
		\indent 在一个二维平面中，有$n$个点。选择其中$m$个点，其中$m$为偶数。对$m$个点编号为$1,2,\cdots,m$，$m$个点的坐标分别为$ (x_{1},y_{1}) , (x_{2},y_{2}) , \cdots , (x_{m},y_{m}) $。将编号相邻的点连接起来，得到奋斗曲线。对于$(x_{i}, y_{i})$ $(2 \le i \le m)$ ，满足以下条件：\\
\begin{equation}\nonumber
	\left\{
	\begin{aligned}
		x_{i-1}<x_{i} \\
		\left\{
		\begin{aligned}
			&y_{i-1}<y_{i} \qquad i \ is \ even \\
			&y_{i-1}>y_{i} \qquad i \ is \ odd
		\end{aligned}
		\right.
	\end{aligned}
	\right.
\end{equation}

		\section*{输入描述}
			\noindent 第1行输入点的数目$n$($n$为整数, $1 \le n \le 1000$)。\\
			第2行-第$n+1$行输入$n$个点的坐标，每行两个整数x,y($0 \le x,y \le 100000$)，分别为各个点的横坐标，纵坐标。保证数据中所有点的坐标都不相同。

		\section*{输出描述}
			\noindent 输出一个整数，为奋斗折线中最大的$m$值。若不存在奋斗折线，请输出$0$。
			
		\section*{示例1}
			\subsection*{输入}
\noindent
5		\\
1 5	\\
2 4	\\
3 3	\\
4 2	\\
5 1

			\subsection*{输出}
\noindent
0

			\subsection*{说明}
			\noindent 不存在奋斗折线

		\section*{示例2}
			\subsection*{输入}
\noindent
7		\\
0 3	\\
3 1	\\
4 2	\\
2 3	\\
3 2	\\
5 -1	\\
1 1


			\subsection*{输出}
\noindent
4

			\subsection*{说明}
			\includegraphics[scale=0.5]{1.png} \\
			\noindent $(1,1),(2,3),(3,1),(4,2)$为m值最大的可行解。\\
			\noindent 另外的可行解有$(1,1),(2,3)$和$(3,1),(4,2)$和$(1,1),(4,2)$。


		\section*{备注}
		\indent 题目数据已降低难度，原来数据$1 \le n \le 100000$ ，有兴趣的同学不妨研究一下。

	\end
	{document}